\section{Related Work}\label{related-work}

Previous research in this area for determining streamflow has used
several different models and traditionally they have been process based
models. These models are typically based on mathematical formulas that
are common in hydrological sciences but there is room for improvement as
they are by no means 100\% accurate. In our main research paper
reference (Kratzert et al., 2019) it is pointed out that machine
learning methods have come a long way and being data driven they take
into account more than the process based methods can. This is
particularly useful for hydrologically diverse catchments where process
based models typically need a deeper knowledge of a particular area
which can be expensive/time consuming.

In the past Machine Learning models have been used to measure streamflow
in ungauged basins even better than conceptual models that were
calibrated for each catchment. This is even against gauged basins that
are operating on a calibrated process model. While in the paper we
researched this was on an average basis (didn't outperform all
process/calibrated models), it does provide evidence that data based
models are effective. This is particularly important for ungauged basins
since even if financially hydrologists are able to start gauging the
basin it will take 2-3 years to get a reliable data set.

Previous machine learning models researched were a type of recurrent
neural network (RNN). They were specifically called Long short-term
memory (LSTM) networks. These were used because they are great for long
term models that look at weather effects that happen over time. They are
useful for things like snowfall accumulation and vegetation growth.
These happen over long lengths of time so it makes them useful for
modeling basin catchment. In the results of their study they observed
that LSTM network models using static catchment attributes was
determined to be the most useful. We will be using this add a reference
to see if we can create a machine learning model that is easier to
interpret as well as less intensive.
